\clearpage\section{附录A：10篇团队论文的完整阅读与迁移索引}
已逐页读取全部10篇PDF，共123页，范围含方法、公式、实验、图表、结论与参考文献；T-ASE另含附录。完整学习笔记见随稿TEAM\_WRITING\_AUDIT.md。此处只列写作迁移对象，不将原文算法或成绩作为V8定位证据。
\begingroup\fontsize{9}{11.5}\selectfont\setlength{\parskip}{1pt}\setlength{\parindent}{0pt}
\begin{longtable}{@{}>{\raggedright\arraybackslash}p{1cm}>{\raggedright\arraybackslash}p{7.0cm}>{\raggedright\arraybackslash}p{7.4cm}@{}}
\toprule
编号/页数 & 原件题录与DOI & 实际阅读重点及限制 \\
\midrule
\endfirsthead
\toprule
编号/页数 & 原件题录与DOI & 实际阅读重点及限制 \\
\midrule
\endhead
\bottomrule
\endfoot

P01\par 6页 & Unmanned Aerial Vehicle Path Planning Based on DP-DDPG Algorithm\par 2023 IEEE International Conference on Unmanned Systems (ICUS), 2023. \par DOI: \url{10.1109/ICUS58632.2023.10318428} & 状态/动作/约束先定义，再奖励、双经验池与噪声修改，最后比较任务成功率与训练收敛。\par 独立模块消融不足；摘要约30\%与表I 51\%→83.25\%应区分百分点及相对改善。 \\\addlinespace[5pt]
P02\par 5页 & Autonomous Unmanned Systems Test Platform for Complex Water Environments\par Guidance, Navigation and Control, 2025. \par DOI: \url{10.1142/S2737480725500359} & 平台功能→硬件/软件→应用验证；明确系统集成和算法本身的证据边界。\par 适合平台描述，不能直接充当三个算法创新的证明模板。 \\\addlinespace[5pt]
P03\par 14页 & A nonlinear model predictive control based control method to quadrotor landing on moving platform\par Cognitive Computation and Systems, 2023. \par DOI: \url{10.1049/ccs2.12081} & 定位、规划、控制三层逐项有模型、有操作、有模块实验，并以完整系统验证收束。\par 部分对既有算法适用范围的概括过强，写作结构可学，技术定性仍须独立核查。 \\\addlinespace[5pt]
P04\par 15页 & 基于多关键点检测加权融合的无人机相对位姿估计算法\par 自动化学报 / Acta Automatica Sinica, 2024. \par DOI: \url{10.16383/j.aas.c230297} & 原流程逐环节缺点图→具体失败→对应操作→五组实验；是本稿第二章最直接的结构参照。\par 去模糊实验使用已知模糊方向/距离等条件，学习时应同时学习暴露前提，不能照搬性能结论。 \\\addlinespace[5pt]
P05\par 13页 & Airborne Camera Dynamic Target Detection Based on Background Prediction and Semantic Compensation in Surface Environment\par IEEE Transactions on Circuits and Systems for Video Technology, 2025. \par DOI: \url{10.1109/TCSVT.2025.3579410} & 将相机自运动造成伪运动和语义类别无法说明瞬时运动分开；输出接口落实到动态特征点剔除/SLAM。\par p.10983有低动态序列收益小甚至下降的边界，不能只学习正结果；个别公式/伪代码记号仍须复核。 \\\addlinespace[5pt]
P06\par 17页 & MARFPNet: Multiattention and Adaptive Reparameterized Feature Pyramid Network for Small Target Detection on Water Surfaces\par IEEE Transactions on Instrumentation and Measurement, 2024. \par DOI: \url{10.1109/TIM.2024.3485463} & 提取/融合/重采样/损失分别对准小目标信息问题，独立模块与组合消融并列，准确率/复杂度同时报告。\par DySample和NWD是有来源的引入；表VIII大中目标上基线更优；同设置重跑与文献结果须区分。 \\\addlinespace[5pt]
P07\par 15页 & A Strong UAV Vision Tracker Based on Deep Broad Learning System and Correlation Filter\par IEEE Transactions on Automation Science and Engineering, 2025. \par DOI: \url{10.1109/TASE.2024.3429161} & 先给DSST实际滤波计算，再给可观测异常量RSV和恢复接口，实验同时检验初始化、状态检测、恢复及最终跟踪。\par p.5721声明DBLS验证集与训练集相同；摘要/实践说明/实验段对比数量口径不同，不应照学。正式卷年2025、online2024不能混淆。 \\\addlinespace[5pt]
P08\par 14页 & A Non-Parametric Scheme for Identifying Data Characteristic Based on Curve Similarity Matching\par IEEE/CAA Journal of Automatica Sinica, 2024. \par DOI: \url{10.1109/JAS.2024.124359} & 样本→密度估计→形状/突变/对称指标→评估，指标及图中中间量前后一致；小样本与在线滑窗区别。\par p值并非样本来自高斯分布的后验概率；直接把CS与不同检验的p值数值比较不能替代统一误报率/检验力验证。 \\\addlinespace[5pt]
P09\par 12页 & 基于PCA和ICA模式融合的非高斯特征检测识别\par 自动化学报 / Acta Automatica Sinica, 2024. \par DOI: \url{10.16383/j.aas.c230326} & 预处理、混合核/优化、ICA-PCA相关性降维、统计检测各有中间指标，再看USV实际观测。\par 保留方差/降维指标不等价于最终任务识别信息；未拒绝正态性不能写成证明正态；技术泛化论断须独立审查。 \\\addlinespace[5pt]
P10\par 12页 & IMM-AIS-UKF: A GAN-Enhanced Dynamics-Constrained Trajectory Predictor With DDPG Method for UAV-USV Interaction\par IEEE Transactions on Intelligent Transportation Systems, 2025. \par DOI: \url{10.1109/TITS.2025.3591077} & 长预测时域转弯偏差→子步积分→精度/计算成本权衡→自适应子步与模型融合；离线训练和在线推理分开。\par 原PDF标future issue/分页未定，最终卷期未确认；GAN完整网络、损失、训练和独立消融在上传全文中未充分展开，不能照搬标题式创新包装。 \\\addlinespace[5pt]
\end{longtable}
\endgroup

\clearpage\section{附录B：正文29项参考文献逐条核验}
正文包含24项学术来源、4项冻结项目证据及1项官方作者指南。元数据确认、方法阅读与复现实验分级记录；没有读到的最终全文不标为全文阅读。完整字段及一级来源URL见CURRENT\_29\_REFERENCE\_REVIEW.csv。
\begingroup\fontsize{9}{11.5}\selectfont\setlength{\parskip}{1pt}\setlength{\parindent}{0pt}
\begin{longtable}{@{}>{\raggedright\arraybackslash}p{0.7cm}>{\raggedright\arraybackslash}p{14.8cm}@{}}
\toprule
序号 & 题录、实际核验与论断范围 \\
\midrule
\endfirsthead
\toprule
序号 & 题录、实际核验与论断范围 \\
\midrule
\endhead
\bottomrule
\endfoot

1 & \textbf{Gas source localization and mapping with mobile robots: A review}\par Francis, Adam and Li, Shuai and Griffiths, Christian and Sienz, Johann. 2022. DOI: \url{10.1002/rob.22109}. \par\textbf{类别/核验：} 同行评议综述；L1；2022；Journal of Field Robotics 39(8):1341–1373；DOI 10.1002/rob.22109。\par\textbf{引用位置：} 第1节“引言”\par\textbf{实际阅读：} 只核Crossref登记的作者、题名、期刊、卷期、页码和DOI；未阅读全文。\par\textbf{论据范围：} 作为移动机器人气体建图与源定位领域的综述入口；本稿对任务估计对象的区分还依赖问题定义本身。\par\textbf{限制/待核：} 未逐段核综述分类和具体结论；不得借此引用具体方法优劣、性能数字或声称已通读。\par\textbf{一级来源：} \href{https://doi.org/10.1002/rob.22109}{来源1} \\\addlinespace[6pt]
2 & \textbf{Information-Driven Gas Distribution Mapping for Autonomous Mobile Robots}\par Gongora, Andres and Monroy, Javier and Rahbar, Faezeh and Ercolani, Chiara and Gonzalez-Jimenez, Javier and Martinoli, Alcherio. 2023. DOI: \url{10.3390/s23125387}. \par\textbf{类别/核验：} 同行评议论文；L1+L2；2023；Sensors 23(12)，文章号5387；DOI 10.3390/s23125387；作者顺序已核。\par\textbf{引用位置：} 第1节“引言”\par\textbf{实际阅读：} Crossref完整元数据与作者/机构出版记录；出版方全文请求429限流，未阅读全文。\par\textbf{论据范围：} 存在面向自主移动机器人的信息驱动气体分布建图研究；与源位置估计的研究对象区分。\par\textbf{限制/待核：} 具体目标函数、信息准则、实现和性能未全文核查。429表示访问限流，不表示DOI无效。\par\textbf{一级来源：} \href{https://www.mdpi.com/1424-8220/23/12/5387}{来源1}；\href{https://www.epfl.ch/labs/disal/research/gassensingstructure/}{来源2} \\\addlinespace[6pt]
3 & \textbf{Application of the NCAR FastEddy Microscale Model to a Lake Breeze Front}\par Welch, Brittany M. and Horel, John D. and Sauer, Jeremy A.. 2024. DOI: \url{10.3390/atmos15070809}. \par\textbf{类别/核验：} 同行评议论文；L1+L3；2024；Atmosphere 15(7)，文章号809；DOI 10.3390/atmos15070809。\par\textbf{引用位置：} 第1节“引言”\par\textbf{实际阅读：} 读出版方HTML摘要、模型与场景相关正文；核Crossref身份；未声称逐页通读全部实验。\par\textbf{论据范围：} 湖风锋中的微尺度风场与局地变化为湖岸气体传播任务提供物理背景。\par\textbf{限制/待核：} 不能据此证明枫树岭S2热力场已资格化、已实现真实湖风，或本方法已在湖岸验证。\par\textbf{一级来源：} \href{https://www.mdpi.com/2073-4433/15/7/809}{来源1} \\\addlinespace[6pt]
4 & \textbf{Measurement report: Observations of ground-level ozone concentration gradients perpendicular to the Lake Ontario shoreline}\par Huang, Yao Yan and Donaldson, D. James. 2024. DOI: \url{10.5194/acp-24-2387-2024}. \par\textbf{类别/核验：} 同行评议论文；L1+L3；2024；Atmospheric Chemistry and Physics 24(4):2387–2398；DOI 10.5194/acp-24-2387-2024。\par\textbf{引用位置：} 第1节“引言”\par\textbf{实际阅读：} 读官方摘要、沿岸采样方法、梯度机制讨论及元数据。\par\textbf{论据范围：} Lake Ontario岸线垂直方向的近地面臭氧浓度梯度和局地气象差异可作为应用背景。\par\textbf{限制/待核：} 臭氧受化学反应与沉降影响；不能等同被动气体单源泄漏实验，也不能外推S2场景验证结果。\par\textbf{一级来源：} \href{https://acp.copernicus.org/articles/24/2387/2024/acp-24-2387-2024.html}{来源1} \\\addlinespace[6pt]
5 & \textbf{3D Gas Sensing with Multiple Nano Aerial Vehicles: Interference Analysis, Algorithms and Experimental Validation}\par Ercolani, Chiara and Jin, Wanting and Martinoli, Alcherio. 2023. DOI: \url{10.3390/s23208512}. \par\textbf{类别/核验：} 同行评议论文；L1+L2；2023；Sensors 23(20)，文章号8512；DOI 10.3390/s23208512。\par\textbf{引用位置：} 第1节“引言”\par\textbf{实际阅读：} 核Crossref和EPFL作者/机构出版记录；出版方429；未通读方法和实验。\par\textbf{论据范围：} 多架微型无人机气体感知已研究平台干扰；支持将感知配置与干扰列为双机实验需控制的因素。\par\textbf{限制/待核：} 不得引用未核的旋翼干扰大小、推荐间距、传感器型号、布局和性能数字；本轮未阅读全文。\par\textbf{一级来源：} \href{https://www.mdpi.com/1424-8220/23/20/8512}{来源1}；\href{https://www.epfl.ch/labs/disal/research/gassensingstructure/}{来源2} \\\addlinespace[6pt]
6 & \textbf{Incorporating Animal Klinotaxis Behavior and Bilateral Concentration Information: A 3-D Bio-Inspired Odor Source Localization Algorithm and Its Performance Evaluation}\par Lin, Qin and Wu, Shuo and Wang, Hui and Ma, Ke and Jiang, Lelun and Zhang, Jinxiu. 2025. DOI: \url{10.1109/JSEN.2024.3407982}. \par\textbf{类别/核验：} 同行评议论文；L1+L2；正式卷期为2025年IEEE Sensors Journal 25(1):697–708；DOI 10.1109/JSEN.2024.3407982；DOI含2024不等于卷期年。旧键Bilateral2024已改为Lin2025Bilateral。\par\textbf{引用位置：} 第1节“引言”\par\textbf{实际阅读：} 核Crossref和IEEE索引摘要/方法片段；全文PDF访问418，未通读最终PDF。\par\textbf{论据范围：} 动物klinotaxis与双侧浓度信息已用于三维气味源定位，双侧配置本身不能作为本文原创。\par\textbf{限制/待核：} 未复现该算法；不将其性能数字转借为本方法效果；引用内容限已核摘要与研究范围。\par\textbf{一级来源：} \href{https://ieeexplore.ieee.org/document/10551708/}{来源1} \\\addlinespace[6pt]
7 & \textbf{Insect-inspired adaptive behavioral compensation strategy against olfactory sensory deficiency for robotic odor source localization}\par Shigaki, Shunsuke and Yokota, Keisuke and Sekiwa, Ryoko and Kurabayashi, Daisuke and Owaki, Dai. 2026. DOI: \url{10.1038/s44182-026-00080-5}. \par\textbf{类别/核验：} 同行评议论文；L2+L3；2026；npj Robotics 4，文章号12；DOI 10.1038/s44182-026-00080-5；五位作者已核。\par\textbf{引用位置：} 第1节“引言”\par\textbf{实际阅读：} 读Nature官方摘要、引言和嗅觉缺损补偿任务/机器人验证范围；未做代码复现。\par\textbf{论据范围：} 昆虫启发行为补偿与嗅觉缺损条件下的机器人搜索已有先例。\par\textbf{限制/待核：} 不能据此声称本文首次提出双侧嗅觉、感知缺损恢复或等同M3源假设恢复；未复现实验。\par\textbf{一级来源：} \href{https://www.nature.com/articles/s44182-026-00080-5}{来源1} \\\addlinespace[6pt]
8 & \textbf{Deep Learning Based Topography Aware Gas Source Localization with Mobile Robot}\par Tian, Changhao and Wang, Annan and Fan, Han and Wiedemann, Thomas and Luo, Yifei and Yang, Le and Lin, Weisi and Lilienthal, Achim J. and Chen, Xiaodong. 2025. DOI: \url{10.1109/ICRA55743.2025.11128134}. \par\textbf{类别/核验：} 同行评议会议论文；L1+L2；ICRA 2025:4380–4386；DOI 10.1109/ICRA55743.2025.11128134；完整九作者身份由登记核对。\par\textbf{引用位置：} 第1节“引言”\par\textbf{实际阅读：} 主代理核Crossref、IEEE官方摘要以及DLR/TUM作者机构页面；未阅读全文。\par\textbf{论据范围：} 地形感知的学习方法已用于移动机器人气体源定位；可支持相关研究方向存在。\par\textbf{限制/待核：} 训练/测试权限、模型内部结构、泛化条件和增益数值未通读核查；不将“改进表达”写成已核的量化性能结论。\par\textbf{一级来源：} \href{https://ieeexplore.ieee.org/document/11128134/}{来源1} \\\addlinespace[6pt]
9 & \textbf{Towards Efficient Gas Leak Detection in Built Environments: Data-Driven Plume Modeling for Gas Sensing Robots}\par Jin, Wanting and Rahbar, Faezeh and Ercolani, Chiara and Martinoli, Alcherio. 2023. DOI: \url{10.1109/ICRA48891.2023.10160816}. \par\textbf{类别/核验：} 同行评议会议论文；L1+L2；ICRA 2023:7749–7755；DOI 10.1109/ICRA48891.2023.10160816；采用出版登记作者顺序Jin、Rahbar、Ercolani、Martinoli。\par\textbf{引用位置：} 第1节“引言”\par\textbf{实际阅读：} 核Crossref与EPFL作者团队页面；主代理读作者页面；未阅读全文。\par\textbf{论据范围：} 建成环境中面向气体感知机器人的数据驱动羽流建模已有研究。\par\textbf{限制/待核：} 未核完整模型、具体环境和实验数字；不能据作者页认定已通读论文或证明本项目物理模型有效。\par\textbf{一级来源：} \href{https://doi.org/10.1109/ICRA48891.2023.10160816}{来源1}；\href{https://www.epfl.ch/labs/disal/research/gassensingstructure/}{来源2} \\\addlinespace[6pt]
10 & \textbf{Robotic Gas Source Localization With Probabilistic Mapping and Online Dispersion Simulation}\par Ojeda, Pepe and Monroy, Javier and Gonzalez-Jimenez, Javier. 2024. DOI: \url{10.1109/TRO.2024.3426368}. \par\textbf{类别/核验：} 同行评议论文与作者代码；L1+L3+L4；2024；IEEE Transactions on Robotics 40:3551–3564；DOI 10.1109/TRO.2024.3426368；作者稿arXiv:2304.08879v3。\par\textbf{引用位置：} 第1节“引言”；第1节“引言”；第2.2节“原生PMFS计算链及具体研究动机”\par\textbf{实际阅读：} 读作者稿III–V节的计算链、相关方法和局限；核官方PMFS关键源码，未将代码阅读称为复现实验。\par\textbf{论据范围：} 历史命中图与置信度、候选源传播预测、逐格兼容评分、粗细网格、源分布加权预测方差和未观测置信度选点；PMFS已有历史和主动搜索。\par\textbf{限制/待核：} 作者代码4e141e1与项目O24来源a8d5a76分开；不能把整个累计命中图重复当新独立似然，不能将本项目V7负结果归因PMFS。\par\textbf{一级来源：} \href{https://arxiv.org/html/2304.08879v3}{来源1}；\href{https://doi.org/10.1109/TRO.2024.3426368}{来源2}；\href{https://github.com/MAPIRlab/GasSourceLocalization}{来源3} \\\addlinespace[6pt]
11 & \textbf{Deep Probabilistic Indoor Gas Source Localization via Physical Dependency-Guided Sequential Inference}\par Kim, Seunghwan and Kim, Hyungjin and Lee, Junhee and Oh, Hyondong. 2026. \par\textbf{类别/核验：} 作者预印本；L2+L3；arXiv:2608.16221v1，2026-08-17；四位作者已核；正式期刊/会议出版状态未确认。\par\textbf{引用位置：} 第1节“引言”\par\textbf{实际阅读：} 读作者arXiv v1问题定义、风→浓度→源条件推断链、主动策略及实验适用范围。\par\textbf{论据范围：} 连续浓度、风与源的物理依赖顺序概率推断已有相关研究；源导向建模不是本文独有方向。\par\textbf{限制/待核：} 明确标预印本，不写成已同行评议；该文性能不能外推到当前两机或S2任务。\par\textbf{一级来源：} \href{https://arxiv.org/html/2608.16221v1}{来源1} \\\addlinespace[6pt]
12 & \textbf{`Infotaxis' as a strategy for searching without gradients}\par Vergassola, Massimo and Villermaux, Emmanuel and Shraiman, Boris I.. 2007. DOI: \url{10.1038/nature05464}. \par\textbf{类别/核验：} 同行评议论文；L2；2007；Nature 445:406–409；DOI 10.1038/nature05464。\par\textbf{引用位置：} 第1节“引言”\par\textbf{实际阅读：} 核Nature官方摘要及书目；未通读完整论文。\par\textbf{论据范围：} 信息获取/不确定性减少驱动的主动嗅觉搜索已有经典先例。\par\textbf{限制/待核：} 本稿仅作方法谱系引用；实际实现infotaxis基线前仍需核其观测模型、动作和目标函数；不声称全文读过。\par\textbf{一级来源：} \href{https://www.nature.com/articles/nature05464}{来源1} \\\addlinespace[6pt]
13 & \textbf{Sense in Motion with Belief Clustering: Efficient Gas Source Localization with Mobile Robots}\par Jin, Wanting and Martinoli, Alcherio. 2024. DOI: \url{10.1109/ICRA57147.2024.10611440}. \par\textbf{类别/核验：} 同行评议会议论文；L1+L2+局部L3；ICRA 2024:14909–14916；DOI 10.1109/ICRA57147.2024.10611440；作者Jin、Martinoli。\par\textbf{引用位置：} 第1节“引言”\par\textbf{实际阅读：} 核Crossref、作者履历、机构摘要；读作者摘要及可检索方法片段；PDF直取429，未逐页通读。\par\textbf{论据范围：} 运动中连续感知、信念聚类以及羽流/传感器时间特性应进入主动定位比较。\par\textbf{限制/待核：} 不能简化为仅在终点一次EIG；不引用未逐项核对的性能数字；未运行算法。\par\textbf{一级来源：} \href{https://infoscience.epfl.ch/entities/publication/22992bbb-65c6-4b2c-a983-17f4050d44f2}{来源1}；\href{https://graphsearch.epfl.ch/en/publication/22992bbb-65c6-4b2c-a983-17f4050d44f2}{来源2} \\\addlinespace[6pt]
14 & \textbf{Cumulative Informative Path Planning for Efficient Gas Source Localization with Mobile Robots}\par Jin, Wanting and Leroy, Hugo and Bösel-Schmid, Nicolaj and Martinoli, Alcherio. 2025. DOI: \url{10.1109/IROS60139.2025.11246224}. \par\textbf{类别/核验：} 同行评议会议论文；L1+L2+局部L3；IROS 2025:18088–18094；DOI 10.1109/IROS60139.2025.11246224；四位作者已核。\par\textbf{引用位置：} 第1节“引言”\par\textbf{实际阅读：} 核Crossref、EPFL条目和作者正式履历；读作者摘要及作者上传稿可得内容；未声称逐页通读最终PDF。\par\textbf{论据范围：} 气体源定位已有沿路径累积信息规划；比较应包括移动、沿途观测和停留成本。\par\textbf{限制/待核：} 不将路径累计信息称为本文首次；未运行作者算法，不引用其性能作为本文证据。\par\textbf{一级来源：} \href{https://infoscience.epfl.ch/entities/publication/3a6f868c-8d68-4aaa-937e-900941d3c694}{来源1}；\href{https://disalw3.epfl.ch/people/alcherio/files/AM_full_CV.pdf}{来源2}；\href{https://www.epfl.ch/labs/disal/research/gassensingstructure/}{来源3} \\\addlinespace[6pt]
15 & \textbf{Fly navigational responses exploit plume-specific odor motion and gradient cues}\par Brudner, Samuel and Zhou, Baohua and Jayaram, Viraaj and Madeira Santana, Gustavo and Crimaldi, John P. and Clark, Damon A. and Emonet, Thierry. 2026. DOI: \url{10.1073/pnas.2608896123}. \par\textbf{类别/核验：} 同行评议论文与作者代码；L1+L3+L4；2026；PNAS 123(29):e2608896123；DOI 10.1073/pnas.2608896123；七位作者已核。\par\textbf{引用位置：} 第1节“引言”；第3节“面向竞争源判别的双侧镜像嗅觉计算”\par\textbf{实际阅读：} 读PNAS主文相关方法、图3模型、作者figure3\_models.py和复现说明；未逐项核完全部SI，未运行母算法实验。\par\textbf{论据范围：} 可学习双侧因果时间滤波、整流、镜像相减与sigmoid中心线侧别判别；梯度和气味运动互补是母文来源。\par\textbf{限制/待核：} V7固定核库与TRAIN选择不是原可学习MNM复现；母文不验证候选源后验、双无人机或M1定位收益；新意限任务/证据/选择机制改造。\par\textbf{一级来源：} \href{https://www.pnas.org/doi/10.1073/pnas.2608896123}{来源1}；\href{https://github.com/emonetlab/GradientMotionMultiplePlumes}{来源2} \\\addlinespace[6pt]
16 & \textbf{A vector-based strategy for olfactory navigation in Drosophila}\par Siliciano, Andrew F. and Minni, Sun and Morton, Chad and Dowell, Charles K. and Eghbali, Noelle B. and Busch, Silas E. and Rhee, Juliana Y. and Abbott, L. F. and Ruta, Vanessa. 2026. DOI: \url{10.1038/s41586-026-10827-7}. \par\textbf{类别/核验：} 同行评议论文元数据＋作者稿/代码；L1+L3+L4；Nature 657:443–454，2026；DOI 10.1038/s41586-026-10827-7；九位作者与正式题名已核。\par\textbf{引用位置：} 第1节“引言”；第4.1节“M2：动态边界方向记忆与返回区域”\par\textbf{实际阅读：} 核最终Nature书目；读作者2025预印本相关方向记忆方法、官方作者代码及出版方评审材料；最终全文跳转受限，未声称通读最终版。\par\textbf{论据范围：} 气味进入/退出事件关联的方向向量记忆和边界返回策略，可启发M2候选动作生成。\par\textbf{限制/待核：} 作者稿与最终版应区分；不能解释为稳定世界坐标边界图、源概率图或已实现双机编队；向量不能另作独立源似然。\par\textbf{一级来源：} \href{https://www.nature.com/articles/s41586-026-10827-7}{来源1}；\href{https://pmc.ncbi.nlm.nih.gov/articles/PMC11844514/}{来源2}；\href{https://github.com/rutalaboratory/Edge-Tracking-Model}{来源3} \\\addlinespace[6pt]
17 & \textbf{Tradeoffs Among Action Taking Policies Matter in Active Sequential Multi-Hypothesis Testing: The Optimal Error Exponent Region}\par Hsu, Chia-Yu and Wang, I-Hsiang. 2025. DOI: \url{10.1109/TIT.2025.3588527}. \par\textbf{类别/核验：} 同行评议论文＋作者稿；L1+L3；2025；IEEE Transactions on Information Theory 71(9):6546–6565；DOI 10.1109/TIT.2025.3588527；作者稿2405.06554v3。\par\textbf{引用位置：} 第1节“引言”；第4.3节“M3：完整源集上的成本分配与持续覆盖”\par\textbf{实际阅读：} 读作者稿问题定义、动作/资源约束与错误指数框架；未声称逐行验证全部定理证明。\par\textbf{论据范围：} 主动多假设检验中的动作策略、资源分配和竞争假设区分已有理论基础。\par\textbf{限制/待核：} 已知条件模型及渐近条件不能直接转成有限时长、时变羽流UAV的最优性或恢复保证。\par\textbf{一级来源：} \href{https://arxiv.org/html/2405.06554v3}{来源1}；\href{https://cc.ee.ntu.edu.tw/~ihsiangw/NICLab/Publications.html}{来源2} \\\addlinespace[6pt]
18 & \textbf{Active Sequential Hypothesis Testing with Non-Homogeneous Costs}\par Vershinin, George and Cohen, Asaf and Gurewitz, Omer. 2026. DOI: \url{10.1109/ICASSP55912.2026.11463839}. \par\textbf{类别/核验：} 同行评议会议论文＋作者稿；L1+L3；ICASSP 2026:91–95；DOI 10.1109/ICASSP55912.2026.11463839；作者稿arXiv:2509.11632v3。\par\textbf{引用位置：} 第1节“引言”；第4.3节“M3：完整源集上的成本分配与持续覆盖”\par\textbf{实际阅读：} 读作者稿相关问题、非同质成本比率目标、改造Chernoff策略及理论假设；核ICASSP官方记录与Crossref。\par\textbf{论据范围：} 成本份额与动作抽样频率需区分，q(a)正比w(a)/c(a)；资源效率目标是期望信息与期望成本之比。\par\textbf{限制/待核：} 不能将E[D]/E[c]写成E[D/c]；已知模型、可实施动作及渐近假设限制移植；不自动证明M3有限样本最优。\par\textbf{一级来源：} \href{https://arxiv.org/html/2509.11632v3}{来源1}；\href{https://www.cmsworkshops.com/ICASSP2026/view_paper.php?PaperNum=7206&bare=1}{来源2}；\href{https://doi.org/10.1109/ICASSP55912.2026.11463839}{来源3} \\\addlinespace[6pt]
19 & \textbf{Universal inference}\par Wasserman, Larry and Ramdas, Aaditya and Balakrishnan, Sivaraman. 2020. DOI: \url{10.1073/pnas.1922664117}. \par\textbf{类别/核验：} 同行评议论文；L1+L3；2020；PNAS 117(29):16880–16890；DOI 10.1073/pnas.1922664117。\par\textbf{引用位置：} 第1节“引言”；第4.4节“M3：新后缀审计与可恢复候选列表”\par\textbf{实际阅读：} 读通用似然构造和适用条件；未声称逐行复核全部证明。\par\textbf{论据范围：} 通用似然比推断为预测备择与零假设似然比较提供基础；本稿须自行列明顺序化与条件校准假设。\par\textbf{限制/待核：} 不得从该引用直接推出任意失配模型中的e过程有效；新后缀本身不能修复模型误差；审计不能再次乘入气体后验。\par\textbf{一级来源：} \href{https://www.pnas.org/doi/10.1073/pnas.1922664117}{来源1} \\\addlinespace[6pt]
20 & \textbf{Quantum Sequential Universal Hypothesis Testing}\par Zecchin, Matteo and Simeone, Osvaldo and Ramdas, Aaditya. 2025. \par\textbf{类别/核验：} 作者预印本；最终出版元数据待核；L2+L3；实际引用读过的arXiv:2508.21594v1，2025-08-29；三位作者已核；不得虚填2026正式卷页/DOI。\par\textbf{引用位置：} 第1节“引言”；第4.4节“M3：新后缀审计与可恢复候选列表”\par\textbf{实际阅读：} 读arXiv v1的算法、可预测似然、复合零假设与两侧检验条件；作者主页列npj Quantum Information 2026，但最终DOI未独立核得。\par\textbf{论据范围：} 借鉴已读作者稿中的可预测替代密度与复合零假设审计结构，不涉及量子硬件。\par\textbf{限制/待核：} 复合零假设分母应对整段前缀共同参数最大化；需真实条件分布包含、可预测归一化和多次审计错误预算；量子实验不是气体定位证据。\par\textbf{一级来源：} \href{https://arxiv.org/html/2508.21594v1}{来源1}；\href{https://web.stanford.edu/~aramdas/}{来源2} \\\addlinespace[6pt]
21 & \textbf{GADEN: A 3D Gas Dispersion Simulator for Mobile Robot Olfaction in Realistic Environments}\par Monroy, Javier and Hernandez-Bennetts, Victor and Fan, Han and Lilienthal, Achim and Gonzalez-Jimenez, Javier. 2017. DOI: \url{10.3390/s17071479}. \par\textbf{类别/核验：} 同行评议论文＋作者项目记录；L1+L2；2017；Sensors 17(7)，文章号1479；DOI 10.3390/s17071479。\par\textbf{引用位置：} 第2.1节“场景、观测任务与固有挑战”\par\textbf{实际阅读：} 核作者项目书目、机构页面和元数据；MDPI请求429，未阅读全文。\par\textbf{论据范围：} GADEN是移动机器人嗅觉的三维气体扩散仿真平台；支持说明仿真工具来源。\par\textbf{限制/待核：} 不能由平台论文推出理想浓度等于真实PID动态，也不能证明本项目源高、湖岸热场和实飞已验证；本轮未读2017全文。\par\textbf{一级来源：} \href{https://www.mdpi.com/1424-8220/17/7/1479}{来源1}；\href{https://mapir.isa.uma.es/mapirpubsite/index.php/publications/show/362}{来源2}；\href{https://github.com/MAPIRlab/gaden}{来源3} \\\addlinespace[6pt]
22 & \textbf{Gaden-RT: A real time and interactive gas dispersion simulator for mobile robotics}\par Ojeda, Pepe and Monroy, Javier and Gonzalez-Jimenez, Javier. 2025. DOI: \url{10.1016/j.softx.2025.102388}. \par\textbf{类别/核验：} 同行评议软件论文＋作者代码文档；L1+L2+L4；2025；SoftwareX 32，文章号102388；DOI 10.1016/j.softx.2025.102388。\par\textbf{引用位置：} 第2.1节“场景、观测任务与固有挑战”\par\textbf{实际阅读：} 核官方书目、作者README和changelog；出版方全文403；未通读论文全文。\par\textbf{论据范围：} 实时交互气体仿真平台来源；作者文档中的可选GPU浓度图加速仅按对应环节描述。\par\textbf{限制/待核：} 不得写成全管线GPU物理求解，也不把代码工具来源当本项目物理资格与定位增益；403不等于文献无效。\par\textbf{一级来源：} \href{https://www.sciencedirect.com/science/article/pii/S2352711025003541}{来源1}；\href{https://github.com/MAPIRlab/gaden}{来源2}；\href{https://github.com/MAPIRlab/gaden/blob/humble/Changelog.md}{来源3} \\\addlinespace[6pt]
23 & \textbf{Source Localization V7：方法、数据合同与实现快照}\par 项目研究记录. 2026. \par\textbf{类别/核验：} 冻结项目代码/研究记录；非同行评议；P；2026项目研究快照；固定提交061ecf5ddea0a59dcf932b5038ad646b7af454a4；不是期刊/会议论文，没有虚构DOI。\par\textbf{引用位置：} 第2.2节“原生PMFS计算链及具体研究动机”；第2.2节“原生PMFS计算链及具体研究动机”；第3.1节“双侧读出、真实时钟与因果镜像算子”；第4节“判别式双机测量与持续反证定位”；第5.2节“固定观测下的完整普通对照”\par\textbf{实际阅读：} 团队审计读取冻结方法文档、代码、合同与结果汇总；文献子任务读跨领域及理论文档；逐文件身份见代码证据清单。\par\textbf{论据范围：} 用于实现状态、固定核库、代码来源、研究原型、已有static pilot与解析矩阵，以及R3汇总的来源和未恢复边界。\par\textbf{限制/待核：} R3窗口级收据未完全恢复；40快照来自4开发世界，不等于40独立世界；156解析episode不是原生物理闭环；设计与软件测试不证明性能。\par\textbf{一级来源：} \href{https://github.com/kris-yun/uav-gsl-isj/tree/061ecf5ddea0a59dcf932b5038ad646b7af454a4/research/source_localization_v7}{来源1} \\\addlinespace[6pt]
24 & \textbf{B4官方终止判定及M0离线归因记录}\par 项目研究记录. 2026. \par\textbf{类别/核验：} 冻结项目失败案例记录；非同行评议；P；B4终止记录提交c2aacbac7fb2ef7be4be3c9c7dcfce20c11f823c；M0归因记录17c48958f9d7eda98dd61b8e656a8c52205cb7d7；2026。\par\textbf{引用位置：} 第2.2节“原生PMFS计算链及具体研究动机”\par\textbf{实际阅读：} 算法审计读取两份冻结原文并核字节SHA；记录来源与当前主张逐项对应，未重跑该实验。\par\textbf{论据范围：} 特定归档运行的success=true、源图方差约0.329、Top-5\%误差5.112m与错/真代表源兼容比约14874.6；集中度与准确性需分开。\par\textbf{限制/待核：} 项目案例不能证明通用统一根因或所有PMFS条件下失败；没有同行评议身份，未新增运行。\par\textbf{一级来源：} \href{https://github.com/kris-yun/uav-gsl-isj/blob/c2aacbac7fb2ef7be4be3c9c7dcfce20c11f823c/research/pmfs_reproduction/b4_official_terminal_resumed_20261010/B4_OFFICIAL_TERMINAL_DECISION_zh.md}{来源1}；\href{https://github.com/kris-yun/uav-gsl-isj/blob/17c48958f9d7eda98dd61b8e656a8c52205cb7d7/research/pmfs_reproduction/b4_m0_offline_attribution_20261010/B4_M0_EVIDENCE_ATTRIBUTION_zh.md}{来源2} \\\addlinespace[6pt]
25 & \textbf{M4传播响应审计与RK2有限响应原型报告}\par 项目研究记录. 2026. \par\textbf{类别/核验：} 冻结项目响应审计记录；非同行评议；P；M4提交5f89ae05e29714cc31d006638c5067e6813e7077；RK2提交5df5771ca916f209d041e3e3ecd548ec65eb8424；2026项目记录。\par\textbf{引用位置：} 第2.2节“原生PMFS计算链及具体研究动机”\par\textbf{实际阅读：} 算法审计读取M4和RK2冻结报告并核字节SHA；按原报告条件引用，不称为本轮新实验。\par\textbf{论据范围：} 限定M4/RK归档：场MSE下降不能自动等同源排名收益；RK2-L相对S连续MSE约降29.22\%、源NLL约1.22\%、Top-1均4/4。\par\textbf{限制/待核：} 不同报告的数据条件与对照须区分；复合分数不自动成为校准后验；不是独立同行评议研究，也非M1增益证据。\par\textbf{一级来源：} \href{https://github.com/kris-yun/uav-gsl-isj/blob/5f89ae05e29714cc31d006638c5067e6813e7077/evidence/causal_compositional_plume_world_model_v1/M4_V2_FROZEN_TRANSPORT_RESPONSE_AUDIT_20260923.md}{来源1}；\href{https://github.com/kris-yun/uav-gsl-isj/blob/5df5771ca916f209d041e3e3ecd548ec65eb8424/research/pmfs_reproduction/PMFS_RK2_FINITE_RESPONSE_PROTOTYPE_20261010/RK2_FINITE_PROTOTYPE_REPORT_zh.md}{来源2} \\\addlinespace[6pt]
26 & \textbf{原生执行报告及冻结R2开发对照收据}\par 项目研究记录. 2026. \par\textbf{类别/核验：} 冻结项目原生实验/收据；非同行评议；P；固定提交061ecf5ddea0a59dcf932b5038ad646b7af454a4；2026-10-23报告及原生R2开发记录。\par\textbf{引用位置：} 第2.2节“原生PMFS计算链及具体研究动机”；第3.3节“联合源推断、信息边界与已发现的退化”；第4节“判别式双机测量与持续反证定位”；第5节“预实验与后续验证方案”\par\textbf{实际阅读：} 团队读冻结执行报告、R2\_RESULT、模型/训练冻结、通道收据与SHA清单；主代理复算已有统计量，未新生成观测。\par\textbf{论据范围：} 8原生世界、14424帧、71软件测试、1200读出回归；两已标定源4DEV理想粗采样下普通单点/完整双点均4/4、M1为2/4、平均MAP误差约2.427m；覆盖阈值饱和、动态通道0。\par\textbf{限制/待核：} 不是PMFS正式闭环对照；M2/M3原生闭环未启动，S2多源DAY/NIGHT未完成；不能推广未知源/风/真实传感器，更不能声称已有定位增益。\par\textbf{一级来源：} \href{https://github.com/kris-yun/uav-gsl-isj/blob/061ecf5ddea0a59dcf932b5038ad646b7af454a4/research/source_localization_v7/NATIVE_EXECUTION_REPORT_zh.md}{来源1} \\\addlinespace[6pt]
27 & \textbf{Elements of Information Theory}\par Cover, Thomas M. and Thomas, Joy A.. 2006. \par\textbf{类别/核验：} 信息论教材；L2；Thomas M. Cover、Joy A. Thomas；Elements of Information Theory，第2版，Wiley，2006，ISBN9780471241959。官方在线首次发布日期2005与纸本版权/产品年2006不同，保留纸本2006。\par\textbf{引用位置：} 第2.3节“科学问题与总体方案”\par\textbf{实际阅读：} 核Wiley官方第二版书目、ISBN和第2章目录含数据处理不等式；未读取付费章全文；本稿短证明属一般推导。\par\textbf{论据范围：} 数据处理不等式的一般出处；固定或仅凭已见训练/历史选定的确定编码不能增加原始观测中关于源的信息。\par\textbf{限制/待核：} 条件历史、动作和固定模型后才应用Markov链；该定理不证明有限模型性能一定不变，也不否认M2新测点可获取新信息；不声称读完整本。\par\textbf{一级来源：} \href{https://onlinelibrary.wiley.com/doi/book/10.1002/047174882X}{来源1}；\href{https://onlinelibrary.wiley.com/doi/abs/10.1002/047174882X.ch2}{来源2} \\\addlinespace[6pt]
28 & \textbf{VGR Dataset: A CFD-based Gas Dispersion Dataset for Mobile Robotic Olfaction}\par Ojeda, Pepe and Monroy, Javier and Gonzalez-Jimenez, Javier. 2023. DOI: \url{10.1007/s10846-023-02012-z}. \par\textbf{类别/核验：} 同行评议数据集论文；L1+L3；2023；Journal of Intelligent \& Robotic Systems 109，文章号75；DOI 10.1007/s10846-023-02012-z；75不是页码75–83。\par\textbf{引用位置：} 第5.1节“数据来源、划分和读出资格”\par\textbf{实际阅读：} 读出版方任务、环境与数据结构相关内容和元数据；不把论文数据集描述当本项目归档结果。\par\textbf{论据范围：} 说明House/VGR/GADEN气体扩散数据背景；本项目8世界、源坐标和帧数由冻结收据单独支撑。\par\textbf{限制/待核：} 不能据VGR论文证明枫树岭S2热力条件、湖岸资格、真实传感器或当前M1增益。\par\textbf{一级来源：} \href{https://link.springer.com/article/10.1007/s10846-023-02012-z}{来源1} \\\addlinespace[6pt]
29 & \textbf{投稿须知}\par 《自动化学报》编辑部. 2026. \par\textbf{类别/核验：} 期刊官方作者指南；非研究论文；G；《自动化学报》编辑部《投稿须知》；2026-10-23为核对日期，不将其伪作页面初次发表年。\par\textbf{引用位置：} 第6节“结论”\par\textbf{实际阅读：} 主代理实际读取2026-10-23可见的官方投稿须知；属于动态政策页面，不是同行评议研究或项目代码。\par\textbf{论据范围：} 投稿格式、双盲、论据与引文可靠性、AI辅助披露及核心内容生成限制等作者要求的现行依据。\par\textbf{限制/待核：} 采用aas格式不表示已符合全部投稿条款，也不表示目标期刊已定；缺失导师“小论文撰写要求PDF”不能由此替代或声称已读。\par\textbf{一级来源：} \href{https://www.aas.net.cn/news/tgxz.htm}{来源1}；\href{https://www.aas.net.cn/news/lunwenmoban.htm}{来源2} \\\addlinespace[6pt]
\end{longtable}
\endgroup

\clearpage\section{附录C：原稿35项参考文献的逐条处置}
以下编号为原稿引用键，保留旧身份以便回查。L1/L2等等级的具体含义见原始核验表。429只表示限流，404也不足以证明文献虚构。删去的支线文献不因此被判断为错误论文；本轮未全文阅读者均明确记录。正文最终是否保留以当前29项表为准。
\begingroup\fontsize{9}{11.5}\selectfont\setlength{\parskip}{1pt}\setlength{\parindent}{0pt}
\begin{longtable}{@{}>{\raggedright\arraybackslash}p{0.7cm}>{\raggedright\arraybackslash}p{7.0cm}>{\raggedright\arraybackslash}p{7.8cm}@{}}
\toprule
旧号 & 原始身份及核验层级 & 纠错、处置与未核事项 \\
\midrule
\endfirsthead
\toprule
旧号 & 原始身份及核验层级 & 纠错、处置与未核事项 \\
\midrule
\endhead
\bottomrule
\endfoot

1 & Francis A, Li S, Griffiths C, Sienz J. Gas source localization and mapping with mobile robots: A review. Journal of Field Robotics, 2022, 39(8): 1341–1373. DOI: 10.1002/rob.22109\par\textbf{核验层级：}L1\par 本轮只核Crossref题名、作者、期刊卷期页码；未阅读全文。 & \textbf{保留} 领域综述入口，用于区分气体分布建图与源定位；具体算法论断应引用原研究。\par 题名、作者缩写、2022年、39(8):1341–1373及DOI一致。\par\textbf{未核：}如引用综述中的具体分类或结论，仍需核相应正文段落。 \\\addlinespace[6pt]
2 & Ercolani C, Jin W, Martinoli A. 3D gas sensing with multiple nano aerial vehicles: Interference analysis, algorithms and experimental validation. Sensors, 2023, 23(20): 8512. DOI: 10.3390/s23208512\par\textbf{核验层级：}L1+L2\par Crossref完整身份；EPFL作者项目页出版记录；未通读该文方法及实验。 & \textbf{保留} 双无人机平台相关背景，提醒旋翼／传感器互扰及真实读出条件需要控制。\par 2023年23(20):8512为文章号；DOI一致。\par\textbf{未核：}互扰大小、感知高度、设备配置和性能数值未核，不能直接写入本稿。 \\\addlinespace[6pt]
3 & Ojeda P, Monroy J, Gonzalez-Jimenez J. Robotic gas source localization with probabilistic mapping and online dispersion simulation. IEEE Transactions on Robotics, 2024, 40: 3551–3564. DOI: 10.1109/TRO.2024.3426368\par\textbf{核验层级：}L3+L4\par 已读作者接收稿III–V节计算链、方法与局限；已读官方PMFS关键源码。 & \textbf{保留（核心）} 母基线的原始计算链及公平比较必须回到本论文和冻结代码。\par 2024年40:3551–3564与DOI一致；纠正正文误述：PMFS已有历史命中图和主动方差选点。\par\textbf{未核：}用户实际PMFS基线需冻结提交；当前humble代码不能替代固定版本。 \\\addlinespace[6pt]
4 & Tian C, Wang A, Fan H, Wiedemann T, Luo Y, Yang L, Lin W, Lilienthal A J, Chen X. Deep learning based topography aware gas source localization with mobile robot. In: 2025 IEEE International Conference on Robotics and Automation. IEEE, 2025: 4380–4386. DOI: 10.1109/ICRA55743.2025.11128134\par\textbf{核验层级：}L1\par 仅Crossref正式登记元数据，未读方法全文。 & \textbf{保留（内容引述待核）} 已有拓扑／地形感知学习型定位的相关工作，可用于避免把学习源位置当作新概念。\par 9作者及ICRA2025:4380–4386、DOI与原稿一致；题名大小写规范化。\par\textbf{未核：}网络结构、数据权限、泛化条件和性能主张尚未核验。 \\\addlinespace[6pt]
5 & Jin W, Martinoli A. Sense in motion with belief clustering: Efficient gas source localization with mobile robots. In: 2024 IEEE International Conference on Robotics and Automation. IEEE, 2024: 14909–14916. DOI: 10.1109/ICRA57147.2024.10611440\par\textbf{核验层级：}L2+L3（局部）\par Crossref、EPFL作者条目及履历；已读作者摘要与可检索方法段落，PDF直取限流，未逐页通读。 & \textbf{保留} 连续运动感知和信念聚类已有先行研究；支撑采样时钟、传感器动态与真实移动预算的讨论。\par ICRA2024:14909–14916、2作者及DOI一致。\par\textbf{未核：}未完成算法复现；不使用其报告增益作为本文证据。 \\\addlinespace[6pt]
6 & Jin W, Leroy H, Bösel-Schmid N, Martinoli A. Cumulative informative path planning for efficient gas source localization with mobile robots. In: 2025 IEEE/RSJ International Conference on Intelligent Robots and Systems. IEEE, 2025: 18088–18094. DOI: 10.1109/IROS60139.2025.11246224\par\textbf{核验层级：}L1+L2\par 初次Crossref429，补核已成功；EPFL项目页、作者履历和作者摘要核对，未通读最终排版全文。 & \textbf{保留} 沿路径累计信息评价已有先例；M2/M3应比较完整移动和观测成本，不能只评价终点。\par IROS2025:18088–18094、4作者及DOI确认；最初429为访问限流，非DOI无效。\par\textbf{未核：}不声称已实现该路径规划对照；具体算法和增益需另读／复现。 \\\addlinespace[6pt]
7 & Ojeda P, Monroy J, Gonzalez-Jimenez J. Gaden-RT: A real time and interactive gas dispersion simulator for mobile robotics. SoftwareX, 2025, 32: 102388. DOI: 10.1016/j.softx.2025.102388\par\textbf{核验层级：}L1+L2+L4（项目说明）\par 出版方元数据、作者README和Changelog已读；最终期刊全文直取403，未全文阅读。 & \textbf{保留} 原生物理开发与复核平台来源。\par SoftwareX32:102388、2025与DOI一致；GPU能力只据官方文档限定到可选浓度图计算。\par\textbf{未核：}不能将平台引用当作本项目S2资格化或实飞有效性的证据。 \\\addlinespace[6pt]
8 & Gongora A, Monroy J, Rahbar F, Ercolani C, Gonzalez-Jimenez J, Martinoli A. Information-driven gas distribution mapping for autonomous mobile robots. Sensors, 2023, 23(12): 5387. DOI: 10.3390/s23125387\par\textbf{核验层级：}L1\par 仅Crossref作者、题名、卷期和文章号；未读全文。 & \textbf{保留（内容引述待核）} 建图与源定位任务的差别；信息驱动采样不是本稿首次引入。\par Sensors2023,23(12):5387为文章号；6作者及DOI一致。\par\textbf{未核：}具体信息目标、建图后验及性能尚未核验，正文只作限定性相关工作说明。 \\\addlinespace[6pt]
9 & Jin W, Rahbar F, Ercolani C, Martinoli A. Towards efficient gas leak detection in built environments: Data-driven plume modeling for gas sensing robots. In: 2023 IEEE International Conference on Robotics and Automation. IEEE, 2023: 7749–7755. DOI: 10.1109/ICRA48891.2023.10160816\par\textbf{核验层级：}L1+L2\par Crossref及EPFL作者项目目录核验；未读全文。 & \textbf{保留（内容引述待核）} 数据驱动传播模型用于结构化环境GSL已有先例，帮助界定源条件模型来源。\par ICRA2023:7749–7755及DOI一致；作者顺序采用出版方登记Jin–Rahbar–Ercolani–Martinoli。\par\textbf{未核：}作者履历个别条目顺序与出版方登记不同；以原论文首页／出版方为准。模型训练和测试划分待核。 \\\addlinespace[6pt]
10 & Hassan S, Wang L, Mahmud K R. Robotic odor source localization via vision and olfaction fusion navigation algorithm. Sensors, 2024, 24(7): 2309. DOI: 10.3390/s24072309\par\textbf{核验层级：}L1\par 只核Crossref身份，未全文阅读。 & \textbf{删除／移入候选文献库} 本稿没有视觉输入或视觉—嗅觉融合模块，删除是范围取舍，非否认论文有效性。\par Sensors2024,24(7):2309及DOI一致。\par\textbf{未核：}如未来加入视觉语义辅助定位，应重新阅读全文并恢复引用。 \\\addlinespace[6pt]
11 & Fukui C, Uchida T, Koizumi S, Murayama Y, Liu H, Nakata T, Terutsuki D. Advanced bio-hybrid drone for superior odor-source localization: High-precision and extended-range detection capabilities. npj Robotics, 2025, 3: 4. DOI: 10.1038/s44182-025-00020-9\par\textbf{核验层级：}L1\par 只核Crossref正式身份，未读方法和实验。 & \textbf{保留（背景限定）} 说明无人机嗅觉感知已有实体研究；不能据此声称本文传感器实现已完成。\par npj Robotics2025,3,文章号4，Crossref文章号与原稿一致；DOI一致。\par\textbf{未核：}设备响应、有效范围和定位精度未核；仅题名背景引用不承载性能数字。 \\\addlinespace[6pt]
12 & Sloman S J, Bharti A, Martinelli J, Kaski S. Bayesian active learning in the presence of nuisance parameters. In: Proceedings of the Fortieth Conference on Uncertainty in Artificial Intelligence. PMLR, 2024, 244: 3245–3263.\par\textbf{核验层级：}L0\par 原稿转录；本轮未查PMLR页面、未读全文。 & \textbf{删除／移入候选文献库} 当前V8方法没有逐式采用该文的干扰参数主动学习算法；避免作为泛化装饰性引用。\par 未提供DOI不代表论文无DOI；UAI卷244和页码3245–3263未独立核。\par\textbf{未核：}如讨论干扰参数主动学习并保留该文，需核PMLR正式条目和具体理论前提。 \\\addlinespace[6pt]
13 & Huan X, Jagalur J, Marzouk Y M. Optimal experimental design: Formulations and computations. Acta Numerica, 2024, 33: 715–840. DOI: 10.1017/S0962492924000023\par\textbf{核验层级：}L0\par 原稿和Crossref429错误记录；本轮未新增全文或出版方读取。 & \textbf{保留（正式引述前待补核）} 普通贝叶斯实验设计／EIG对照的背景，可连接M2/M3与已有主动设计。\par DOI10.1017/S0962492924000023当前未独立核，不因429判无效。\par\textbf{未核：}作者全名、2024年33:715–840及所引结论均待出版方确认。 \\\addlinespace[6pt]
14 & Welch B M, Horel J D, Sauer J A. Application of the NCAR FastEddy microscale model to a lake breeze front. Atmosphere, 2024, 15(7): 809. DOI: 10.3390/atmos15070809\par\textbf{核验层级：}L1+L3\par MDPI出版方正文已获取并阅读摘要、模型设置及背景；Crossref补核成功。 & \textbf{保留（场景动机）} 湖风锋和建筑湍流耦合说明目标场景需要局地风与热边界条件验证。\par 2024年15(7):809与DOI确认，FastEddy注册商标可在书目按排版需要处理。\par\textbf{未核：}不验证枫树岭S2；不能将大盐湖一次事件外推为本文所有DAY/NIGHT条件。 \\\addlinespace[6pt]
15 & Norzam W A S, Hawari H F, Kamarudin K, Juffry Z H M, Hussein N A A, Gupta M, Abdullah A N. Mobile robot gas source localization using SLAM-GDM with a graphene-based gas sensor. Electronics, 2023, 12(1): 171. DOI: 10.3390/electronics12010171\par\textbf{核验层级：}L1\par Crossref身份核验；未读全文。 & \textbf{删除／移入候选文献库} 本文未采用石墨烯传感器或SLAM-GDM实现，硬件支线不作为三创新依据。\par Crossref在线日期为2022-12-30，而卷12期1对应原稿2023年；不能仅凭上线日强改卷期年，最终引用以出版方卷期页为准。\par\textbf{未核：}若未来恢复引用，需核2023卷期著录与具体硬件动态，不声称年错误已定。 \\\addlinespace[6pt]
16 & Andreou M, Chen N, Bollt E. Assimilative causal inference. Nature Communications, 2026, 17: 1854. DOI: 10.1038/s41467-026-68568-0\par\textbf{核验层级：}L0\par 原稿转录及Crossref429；未读正文。 & \textbf{删除} 旧题名“因果推断”的支撑文献；V8未建立独立因果识别创新，不以因果标签包装源概率推断。\par 2026年Nature Communications17:1854及DOI本轮未核；429不证明不真实。\par\textbf{未核：}此删除属于题目和方法归属变化，不对该论文内容或真实性下结论。 \\\addlinespace[6pt]
17 & 何丽, 杜洋鋆, 李知远, 冉腾, 肖文东, 姚佳程. 移动机器人主动嗅觉气源定位方法研究综述. 计算机工程与应用, 2025, 61(10): 36–49.\par\textbf{核验层级：}L2（作者摘要；出版社待核）\par 授权数据库作者摘要、CNKI记录及其他期刊正式引文页交叉核；未通读全文。 & \textbf{保留（中文综述）} 中文主动嗅觉背景及烟羽发现／跟踪／源确认任务梳理；具体算法证据仍引用原研究。\par 第二作者应为杜洋鋆。候选DOI10.3778/j.issn.1002-8331.2410-0128未获主站核验，安全著录不填DOI。\par\textbf{未核：}期刊主站503/502、Crossref404不构成DOI无效结论；完整终版和DOI印刷待核。 \\\addlinespace[6pt]
18 & 李世钰, 袁杰, 谢霖伟, 郭旭, 张宁宁. 局部风信息启发的AVW-PPO室内气源定位算法. 哈尔滨工业大学学报, 2025, 57(8): 57–68.\par\textbf{核验层级：}L0\par 原稿转录；未检索出版方，未读全文。 & \textbf{删除／移入待核候选库} V8不采用PPO训练；压缩与当前三机制没有直接运算关系的强化学习支线。\par 未提供DOI；作者、2025年57(8):57–68均待正式条目确认。\par\textbf{未核：}不能用未核条目评价PPO或所有主动搜索方法的不足。 \\\addlinespace[6pt]
19 & 许顺, 黄宇鼎, 范靖, 汤达恒, 吕游. 三维环境中无人机气味源定位研究. 电子测量与仪器学报, 2025, 39(7): 236–246. DOI: 10.13382/j.jemi.B2508218\par\textbf{核验层级：}L0\par 原稿与Crossref429；未读全文。 & \textbf{移出正文／保留待核候选} 题目与UAV定位直接相关，值得后续补读；本轮缺一级身份和算法内容核验，暂不承担方法论据。\par DOI10.13382/j.jemi.B2508218尚未核；429仅为限流。\par\textbf{未核：}作者、39(7):236–246和三维定位方式待期刊原文核对。 \\\addlinespace[6pt]
20 & 张学锋, 唐晶晶, 江军, 陈迪. 大中型化工装置区多无人机查找泄漏源仿真研究. 中国安全科学学报, 2026, 36(5): 56–63. DOI: 10.16265/j.cnki.issn1003-3033.2026.05.1091\par\textbf{核验层级：}L0\par 原稿与Crossref404；未查期刊全文。 & \textbf{删除／移入第二篇候选库} 题名指向装置区多机任务；当前第一篇限定两机同步测量，不展开集群管理。范围判断仅据题名。\par Crossref未收录该请求不证明DOI无效；2026年36(5):56–63暂不确认为已核。\par\textbf{未核：}保留原始著录线索，后续多机论文需查出版社／CNKI原文。 \\\addlinespace[6pt]
21 & 史天乐, 李飞, 陈〔原稿稀有字宏raresheng，待核〕, 卢春喜, 王维. 基于CFD模拟的人工神经网络动态溯源模型. 化工进展, 2025, 44(8): 4772–4784. DOI: 10.16085/j.issn.1000-6613.2025-0551\par\textbf{核验层级：}L0\par 原稿和Crossref429；未读全文。 & \textbf{删除／移入候选文献库} 当前不训练CFD监督的ANN动态溯源网络；避免由预测网络背景转移论文重点。范围判断据题名。\par 第三作者原稿使用raresheng稀有字宏，具体字形未核，不能猜字；其余身份及DOI均待核。\par\textbf{未核：}如恢复引用，应先读原首页核作者字形、页码和预测／反演任务。 \\\addlinespace[6pt]
22 & 陈增强, 齐聪聪, 赵怡雯, 程义. 基于多策略改进哈里斯鹰优化算法的危化品泄漏源反演. 上海交通大学学报, 2026, 60(5): 860–870. DOI: 10.16183/j.cnki.jsjtu.2024.452\par\textbf{核验层级：}L0\par 原稿和Crossref429；未读全文。 & \textbf{删除／移入候选文献库} V8没有哈里斯鹰优化模块，题名所示启发式反演不构成本稿三机制直接来源。\par 不能仅凭DOI含2024就否定原稿2026卷期年；正式刊期与早期年份关系待核。\par\textbf{未核：}60(5):860–870、作者及DOI待出版方确认；未做算法优劣判断。 \\\addlinespace[6pt]
23 & 江逸远, 谷小婧, 顾幸生. 基于红外视频的VOCs泄漏源定位与气羽实例分割. 华东理工大学学报(自然科学版), 2024, 50(5): 695–707. DOI: 10.14135/j.cnki.1006-3080.20231011002\par\textbf{核验层级：}L0+L2（页面线索）\par 原稿与Crossref429；检索曾出现该期刊官方页面，但未读该文方法或通篇核验。 & \textbf{删除／移入未来任务候选库} 红外视频和气羽实例分割与当前二维源概率输出不同，不加入未经验证的完整羽流重建主张。\par 已有官方URL线索；本轮未完整核对2024年50(5):695–707及原作者著录。\par\textbf{未核：}未检验红外方法具体贡献，不以题名范围判断代替全文评价。 \\\addlinespace[6pt]
24 & 李鹏, 单钰强. 基于协方差拟合的离格气体泄漏源定位方法. 电子器件, 2023, 46(2): 478–484. DOI: 10.3969/j.issn.1005-9490.2023.02.028\par\textbf{核验层级：}L0\par 原稿与Crossref429；未读全文。 & \textbf{删除／移入候选文献库} 本文未采用协方差拟合和离格阵列反演；删除是篇幅与运算接口的取舍。\par 作者、2023年46(2):478–484与DOI待核；不因429判无效。\par\textbf{未核：}若讨论固定阵列反演，需补查原文和阵列假设。 \\\addlinespace[6pt]
25 & 王蕊, 吴玲玲, 李淑娟, 张欢, 徐智霞. 基于无线传感网络的重气泄漏监测与定位. 火箭军工程大学学报, 2025, 39(2): 70–78. DOI: 10.20189/j.cnki.CN/61-1527/E.202502008\par\textbf{核验层级：}L0\par 原稿与Crossref404；未读全文。 & \textbf{删除／移入候选文献库} 无线固定传感网络和重气扩散不是本稿当前同步UAV测量与已资格化模型。范围判断据题名。\par Crossref404不等于DOI无效；带斜线刊号的DOI原样保留用于后续核对。\par\textbf{未核：}作者、39(2):70–78、重气模型假设待原文确认。 \\\addlinespace[6pt]
26 & 陈泽众, 姚逸卿, 阳媛, 鲍琳欣. 基于时间加权最大似然估计的室内气体源定位方法. 仪器仪表学报, 2025, 46(2): 92–102.\par\textbf{核验层级：}L0\par 原稿转录；未查出版社、未读全文。 & \textbf{移出正文／保留待核候选} 时间加权似然与历史观测相关，但尚未读原算法，不能将其作为已明确超越的对照。\par 未提供DOI不等于无DOI；2025年46(2):92–102及作者待核。\par\textbf{未核：}后续若比较历史加权，应补读其似然、时间权重和观测假设。 \\\addlinespace[6pt]
27 & 董红召, 金灿, 唐伟, 佘翊妮, 林盈盈. 基于遗传-模式搜索算法的微尺度管控区域大气污染物PM\$\_\{2.5\}\$溯源. 浙江大学学报(工学版), 2024, 58(6): 1296–1304. DOI: 10.3785/j.issn.1008-973X.2024.06.019\par\textbf{核验层级：}L0\par 原稿与Crossref429；未读全文。 & \textbf{删除／移入候选文献库} PM2.5区域污染物溯源与局地气体泄漏源辨识的对象及模型不同，当前无直接算法继承。\par 2024年58(6):1296–1304及DOI待核；不把429标为无效。\par\textbf{未核：}删除不评价遗传—模式搜索性能，也不声称已全文审阅。 \\\addlinespace[6pt]
28 & Wang S, Sun S, Liu M, Gao B, Wang Y, Lin F. Resource-aware probability-based collaborative odor source localization using multiple UAVs. China Communications, 2025, 22(12): 269–280. DOI: 10.23919/JCC.ja.2023-0208\par\textbf{核验层级：}L1\par Crossref身份、题名、卷期页码核验；未读方法全文。 & \textbf{保留（多机背景限定）} 资源感知的概率多UAV协作已有先例，双机设备数本身不能成为本文创新。\par 22(12):269–280、2025及DOI一致。登记given/family疑似按中文姓名顺序拆分，不据此自动倒换原稿作者。\par\textbf{未核：}具体作者标准罗马字顺序待原首页；资源定义、协作机制和与两机编队差别待阅读全文。 \\\addlinespace[6pt]
29 & Pearl J. Causality: Models, Reasoning, and Inference. 2nd ed. Cambridge: Cambridge University Press, 2009.\par\textbf{核验层级：}L0\par 本轮仅转录原稿，未新增核出版方版本页。 & \textbf{删除} 不再以候选源因果识别为论文主线，也未使用do算子／因果识别定理，无须保留因果教材作为装饰性基础。\par 第二版2009著录待正式书目核对；不凭一般熟悉程度冒充本轮核验。\par\textbf{未核：}删除只反映V8不采用该理论，非否定因果推断方法。 \\\addlinespace[6pt]
30 & Jeffreys H. Theory of Probability. 3rd ed. Oxford: Clarendon Press, 1961.\par\textbf{核验层级：}L0\par 原稿转录，未查出版社或阅读全文。 & \textbf{删除} 本稿不展开概率哲学或Jeffreys先验；用具体后验、模型与审计原始论文支撑方法。\par 1961第三版、Clarendon Press本轮未核。\par\textbf{未核：}若后续确实采用Jeffreys先验，应查对应章节再引用。 \\\addlinespace[6pt]
31 & Shannon C E. A mathematical theory of communication. Bell System Technical Journal, 1948, 27(3): 379–423.\par\textbf{核验层级：}L0\par Shannon原稿著录未重新核；新替代教材已查Wiley出版社，见本表附录。 & \textbf{替换为Cover与Thomas二版} V8所需是确定性编码的数据处理不等式，使用直接讨论DPI的教材章节更贴近论据。\par 原稿379–423是第一部分的范围线索；若泛指全部论文需另核分期范围，不凭记忆补页。\par\textbf{未核：}替换不表示Shannon论文错误；仅调整论据与引用对应。 \\\addlinespace[6pt]
32 & Vergassola M, Villermaux E, Shraiman B I. ``Infotaxis'' as a strategy for searching without gradients. Nature, 2007, 445(7126): 406–409. DOI: 10.1038/nature05464\par\textbf{核验层级：}L2\par Nature官方摘要及元数据已核；未逐页读原文。 & \textbf{保留} 信息导向气源搜索的经典对照，承认熵下降选点已有来源。\par 2007年445:406–409、DOI一致；原稿7126期号本次未单独查证，可省略期号。\par\textbf{未核：}如复现infotaxis应进一步核观测模型和动作成本，不只引用标题。 \\\addlinespace[6pt]
33 & Thrun S, Burgard W, Fox D. Probabilistic Robotics. Cambridge, MA: MIT Press, 2005.\par\textbf{核验层级：}L0\par 原稿转录；本轮未查MIT Press页面。 & \textbf{删除／备选教材} 研究稿已有后验公式定义和直接PMFS来源，暂不保留未对应具体定理的通用机器人教材。\par 2005年MIT Press与作者本轮未独立核。\par\textbf{未核：}如正文引用标准递归贝叶斯滤波定理，可补查版本和章节后恢复。 \\\addlinespace[6pt]
34 & Brier G W. Verification of forecasts expressed in terms of probability. Monthly Weather Review, 1950, 78(1): 1–3.\par\textbf{核验层级：}L0\par 原稿转录；未查AMS原始出版页。 & \textbf{保留（指标来源待补核）} 若实验方案采用Brier评分，应保留其原始来源而非把该指标称为本文提出。\par 1950年78(1):1–3的著录本轮未独立核；不自造DOI。\par\textbf{未核：}正式定稿前补核原始出版页；多分类Brier所用归一化也须在正文明确。 \\\addlinespace[6pt]
35 & Metzner P, Schütte C, Vanden-Eijnden E. Transition path theory for Markov jump processes. Multiscale Modeling \& Simulation, 2009, 7(3): 1192–1219. DOI: 10.1137/070699500\par\textbf{核验层级：}L0\par 原稿与Crossref429；未阅读全文。 & \textbf{删除} 旧PIP／传播路径方案的理论背景；V8不实现转移路径理论，不应继续暗示该理论支持M3审计。\par 2009年7(3):1192–1219及DOI待核；429不代表无效。\par\textbf{未核：}删除依据当前算法无继承关系，不能写成已读全文后的否定评价。 \\\addlinespace[6pt]
\end{longtable}
\endgroup

\clearpage\section{附录D：待证明命题与禁止升级的表述}
\hypertarget{ux5c1aux5f85ux5b9eux9a8cux8bc1ux660eux7684ux5173ux952eux7ed3ux8bba}{%
\label{ux5c1aux5f85ux5b9eux9a8cux8bc1ux660eux7684ux5173ux952eux7ed3ux8bba}}

本表区分``运算已定义''``条件性质成立''和``性能得到验证''。日期2026-10-23。

\begin{longtable}[]{@{}
  >{\raggedright\arraybackslash}p{(\columnwidth - 8\tabcolsep) * \real{0.2000}}
  >{\raggedright\arraybackslash}p{(\columnwidth - 8\tabcolsep) * \real{0.2000}}
  >{\raggedright\arraybackslash}p{(\columnwidth - 8\tabcolsep) * \real{0.2000}}
  >{\raggedright\arraybackslash}p{(\columnwidth - 8\tabcolsep) * \real{0.2000}}
  >{\raggedright\arraybackslash}p{(\columnwidth - 8\tabcolsep) * \real{0.2000}}@{}}
\toprule\noalign{}
\begin{minipage}[b]{\linewidth}\raggedright
编号
\end{minipage} & \begin{minipage}[b]{\linewidth}\raggedright
需要证明的命题
\end{minipage} & \begin{minipage}[b]{\linewidth}\raggedright
当前证据与状态
\end{minipage} & \begin{minipage}[b]{\linewidth}\raggedright
必要验证
\end{minipage} & \begin{minipage}[b]{\linewidth}\raggedright
目前禁止的表述
\end{minipage} \\
\midrule\noalign{}
\endhead
\bottomrule\noalign{}
\endlastfoot
H1a & 源判别镜像表示较完整raw-joint更可靠 &
固定库与联合表已实现；原生M1仅2/4，普通强对照4/4 &
同观测、同模型头/容量/训练世界、非饱和留出任务 & M1已提高定位精度 \\
H1b & 动态镜像滤波具有独立作用 & 当前选择集合为空，没有动态通道进入评价
& 前瞻冻结新目标；静态/动态/打乱/交换/去共模消融 &
本次实验验证了动态镜像机制 \\
H1c & 源对选择优于母算法直接迁移 & 仅有固定核库best-one，且选择静态梯度
& PNAS可学习结构的公平任务适配及等容量对照 &
已完整复现并超过PNAS母方法 \\
H1d & 有限训练预算下样本效率提高 & 每源TRAIN只有2个世界；无学习曲线 &
世界级独立训练规模、固定最终测试集及区间 &
有统计稳定性或可泛化10\%提升 \\
H1e & 可推广到未见源/风 & 有限表按两个源编号标定 &
参数化源条件预测、未见源/风/realization独立划分 &
已经实现连续源定位泛化 \\
H2a & 方向记忆改善源辨识测量 & 解析原型存在，原生M2闭环未启动 &
同双机/原始流/完整移动停留成本，记忆与无记忆对照 &
原生双机已验证优于普通重捕获 \\
H2b & 双机几何评分对应实际传感器可分性 &
均值/梯度与协方差工程接口；短基线有条件推导 &
真实响应延迟、旋翼、同步、协方差与线性化半径校准 &
长基线由一阶近似无条件认证 \\
H2c & 记忆返回区域具有概率覆盖率 & 只是不确定性传播代理 &
世界级边界事件与覆盖校准 & 返回管道是95\%可信气味边界 \\
H3a & 成本约束反证优于联合EIG & LP等价推导成立，只有解析控制 &
全部合法源、正常/错误先验、同预算和信息权限 & 已证明有限预算全局最优 \\
H3b & 低后验真源能持续恢复 & 有覆盖访问界；当前列表可重新进入 &
可辨动作、正初始质量、付费覆盖、恢复时间/失败率 &
任何错误后验都能恢复；零质量源自动复活 \\
H3c & 新后缀审计控制真实误拒概率 & 条件族内点态支配/Ville界成立 &
真实条件模型包含、时序依赖、支持、重启预算校准 &
拟合直方图即可无条件控制物理风险 \\
Hfull & 完整方法较PMFS-2U有独立收益 & 尚无正式原生比较 &
PMFS-1U/2U、raw-joint+EIG、直接迁移、模块消融 & 相对PMFS已提高10\% \\
Hlake & House结论可迁移到枫树岭S2 & S2完整DAY/NIGHT与多源资格未完成 &
合法源高度/湖陆风/三维读出/独立传播世界资格与迁移 & 已经完成湖岸验证 \\
Hfield & 可实飞并预测完整气团态势 & 本稿无原生飞行或实机数据 &
安全路径、同步到达、传感器校准和任务专属标注 &
已经重建三维羽流或验证实飞性能 \\
\end{longtable}

当前可以保留的结论：确定性编码不增加原始信息；镜像算子满足输入交换反对称；固定总Dirichlet质量的嵌套量化具有一致边缘化；成本比率与对应LP等价；在文中完整假设下，审计量受标准似然比检验过程支配。上述数学结论均不替代性能实验。

\clearpage\section{附录E：编译与交付核查记录}
检查日期：2026年10月23日。以下记录针对最终研究稿，不涉及重新运行V7定位或物理仿真。
\begin{longtable}{@{}>{\raggedright\arraybackslash}p{3.0cm}>{\raggedright\arraybackslash}p{12.5cm}@{}}
\toprule 核查项 & 实际结果 \\\midrule\endhead
编译 & XeTeX 3.141592653--2.6--0.999995（TeX Live 2023/Debian）；原aas类不变，使用匹配版本ctex/xeCJK和Fandol。干净副本连续3次XeLaTeX主文编译均退出0；核验报告连续2次编译退出0。\\
主文结构 & 16页A4；六章、5幅图、3张表、30个编号公式、2段算法流程、29项参考文献。全部43个标签唯一；正文引用均可解析。\\
逐页视觉 & 检查主文全部16页和报告全部页面，并放大核对长公式、表格和参考文献。五图位于第3、5、7、8、10页；三表位于第6、13、14页；主参考文献为第15--16页。没有发现遮挡、越界、缺字或未解析的问号引用。\\
编译诊断 & 最后日志无Overfull盒、Missing character、未定义引用/文献和重复标签。保留的Underfull提示属于双栏/窄表及参考文献疏排；字体Script CJK及形状替代提示已结合视觉检查，不等于缺字。\\
字体 & 主文与报告的字体全部嵌入。英文使用TeX Live字体；中文使用Fandol；核验表作者姓名中的“鋆”由随稿OFL许可的小型重命名字体子集补齐。没有依赖Windows本地字体路径。\\
页眉与元数据 & 主文页眉标为V8研究稿和日期，删除原模板虚拟卷期、作者及DOI字段；PDF作者元数据为空。这不等于已完成公开仓库及全部补充文件的双盲匿名化。\\
源文件可追溯 & 原ZIP未改；原main.tex、原图、原PPTX及章节差异保留。51份V7快照和6份历史报告与远端Git blob一致。14,424帧的SHA核验结论来自冻结记录，本轮没有重跑归档。\\
文献与学术审查 & 10篇团队论文共123页完整阅读；当前29项及旧稿35项逐条列出核验层级、引用用途与未核事项。独立审查的2项阻断和9项重要问题已经逐条落入正文/新图，见随稿最终科学闭合记录。\\
结果边界 & 保留4/4、4/4与2/4负结果以及0动态镜像通道；M2/M3原生闭环、湖岸S2及实飞待执行。数学条件、软件核验与性能结论严格区分。\\
复现入口 & Overleaf上传整包后选择main.tex及XeLaTeX；现成PDF图使正文编译不依赖Python或Office。文献维护脚本与图形TikZ源码、字体许可均随包提供。\\
\bottomrule
\end{longtable}
独立“小论文撰写要求PDF”仍缺失；本文只使用已恢复要求。当前文件用于研究讨论与下一轮实验命题冻结，不是性能已通过、匿名化已完成或AI政策已经审核通过的投稿稿。
